\honest{Science Doesn't Trust Your Rationality}{Eliezer Yudkowsky}{2008}

Scott Aaronson compared the many-worlds interpretation of quantum mechanics with libertarianism: both start from premises most educated people accept, and reach conclusions most of them reject. I quote that much. Aaronson's point, in the same post, was that when very smart people reason their way to such conclusions, you ``would've been wise to doubt'' them even if you could not spot a flaw. I leave that out.

I have argued before that Science rejects many-worlds and Bayes accepts it. ``Science'' has a capital S here: it means an ideal that I describe, not what any scientist actually does. So no scientist can contradict what I say Science says.

My own analogy is between libertarianism and Science. Both, ``it furthermore seems to me,'' distrust reasonable-sounding arguments, try to build systems more trustworthy than their members, and use human flaws to drive the system.

Then a paragraph of politics. When a person or group gets enough power to do everything they think is a good idea, ``history says'' you get Revolutionary France or Soviet Russia. So regulate as little as possible, because you cannot trust lovely theories. ``Regulate as little as possible'' is itself a lovely theory about how much better society would be under one rule. Markets, I add, turn selfishness into production ``With violence restrained and contracts enforced,'' which is the work of courts and police. As for market failures, ``I'm not going to go into'' them.

Science, I say, began as a rebellion against trusting Aristotle, and it lasted because its founders said, ``Let us trust no one! Not even ourselves!'' We do not try to choose impartial judges, I say. Journals appoint editors and referees for exactly that purpose.

Instead, Science harnesses each scientist's stubborn wish to prove his own theory: make a new prediction, do the experiment, and if the result is replicated, you win. This account of science is Michael Polanyi's, from 1962. And a scientist who wants to prove a theory is also motivated to design experiments that the theory will pass, as I admit two paragraphs later.

The system, I say, is ``somewhat more trustworthy than its parts.'' Both systems secretly rely on most people being decent, like tipping at a restaurant you will never visit again.

Then a rule. Successful predictions belong to the theory that made them first, and a later theory that makes no new prediction cannot take them. This is ``an important feature.'' Yesterday I wrote that many-worlds ``doesn't seem to make any new predictions relative to the old theory.'' By this rule, Science is right to withhold its credit.

So Science, I say, ``is not easily reconciled with probability theory.'' A correct probability calculation gives the rational answer, which is true because that is what I mean by rational. Letting Bayes override Science ``is not a trivial step!'' Science assumes that you are ``too stupid and self-deceiving'' to reason perfectly on your own.

I end: ``So, are you going to believe in faster-than-light quantum `collapse' fairies after all? Or do you think you're smarter than that?'' The whole essay has argued that Science is right not to trust how smart you think you are.
